% appendix-numerical-three-valued.tex — \input'd from appendix-numerical.tex
%
% Three-valued predicate relaxation for finite-precision arithmetic.
% Translates each exact Boolean predicate in Algorithm~\ref{alg:ehz}
% into a three-valued version (TRUE / FALSE / UNKNOWN).
%
% Status: draft, not yet Jörn-reviewed.

The exact algorithm (Algorithm~\ref{alg:ehz}) makes several
Boolean decisions that cannot be evaluated exactly in
floating-point arithmetic.
We replace each Boolean predicate~$P$ with a three-valued
predicate~$\tilde P$ taking values in
$\{\textsc{true}, \textsc{false}, \textsc{unknown}\}$,
subject to the \emph{soundness requirement}: whenever $P$ is
true in exact arithmetic, $\tilde P \neq \textsc{false}$.
This ensures that the numerical algorithm searches at least as
broadly as the exact algorithm: wherever the exact algorithm
requires $P$ to hold, the numerical version accepts both
$\textsc{true}$ and $\textsc{unknown}$.
The cost is that $\textsc{unknown}$ may cause the numerical
algorithm to examine pairs $(S, \sigma)$ that the exact
algorithm would have discarded (false positives), but it never
misses a pair that the exact algorithm would have kept.

\paragraph{Adjacency (search-space reduction).}
The adjacency pruning
(Corollary~\ref{cor:adjacency-pruning})
discards pairs where consecutive facets are non-adjacent.
Adjacency is determined combinatorially from the polytope's
vertex--facet incidence: two facets are adjacent if and only if
they share at least one vertex.
Vertex--facet incidence is tested by checking
$\langle n_i, v \rangle = h_i$ (up to a tolerance), so the
three-valued version is:
\begin{align*}
  \tilde P_{\text{adj}}(i,j) &=
  \begin{cases}
    \textsc{true}
      & \text{if some vertex } v \text{ satisfies }
        |\langle n_i, v \rangle - h_i| < \varepsilon
        \text{ and }
        |\langle n_j, v \rangle - h_j| < \varepsilon, \\
    \textsc{false}
      & \text{if } \min_v \max\bigl(
        |\langle n_i, v \rangle - h_i|,\;
        |\langle n_j, v \rangle - h_j|\bigr)
        > \varepsilon', \\
    \textsc{unknown}
      & \text{otherwise (some vertex is near the boundary
        of incidence for at least one facet).}
  \end{cases}
\end{align*}
In practice, for polytopes with vertices of moderate norm
($\|v\| \lesssim 100$) and heights $h_i \in [0.1, 10]$,
the inner products $\langle n_i, v \rangle$ are well-separated
from~$h_i$ for non-incident vertices, so the $\textsc{unknown}$
case rarely arises.

\paragraph{Positivity filter.}
The filter in Step~3 of Algorithm~\ref{alg:ehz} requires
$\beta_i > 0$ for all $i \in S$.
The computed~$\hat\beta$ differs from the true~$\beta^*$ by a
rounding error that depends on the condition number of the
KKT system.
The three-valued version is:
\[
  \tilde P_{>0}(\hat\beta_i) =
  \begin{cases}
    \textsc{true}
      & \text{if } \hat\beta_i > \varepsilon_\beta, \\
    \textsc{false}
      & \text{if } \hat\beta_i < -\varepsilon_\beta, \\
    \textsc{unknown}
      & \text{if } |\hat\beta_i| \leq \varepsilon_\beta,
  \end{cases}
\]
where $\varepsilon_\beta$ is an error bound for the solve.
The $\textsc{unknown}$ case arises when the true~$\beta^*_i$
is near zero:
\begin{itemize}
\item If $0 < \beta^*_i < \varepsilon_\beta$:
  the pair is genuinely needed but may be incorrectly
  discarded ($\hat\beta_i$ could round negative).
\item If $-\varepsilon_\beta < \beta^*_i < 0$:
  the pair is genuinely unnecessary but may be incorrectly
  kept ($\hat\beta_i$ could round positive).
\item If $\beta^*_i = 0$:
  the pair is covered by a smaller
  pair $(S \setminus \{i\}, \sigma|_{S \setminus \{i\}})$,
  which the algorithm also examines.
\end{itemize}
When $\tilde P_{>0}$ returns $\textsc{unknown}$ for some~$i$,
the algorithm keeps the pair but flags it as uncertain
(see Section~\ref{app:error-tracking}).

\paragraph{Near-singular dismissal.}
(Note: unlike the previous two predicates, where
$\textsc{true}$ means ``keep the pair,''
here $\textsc{true}$ means ``safe to dismiss.'')
Lemma~\ref{lem:rank-deficiency-dismissal} allows dismissing
rank-deficient systems, and
Proposition~\ref{prop:approximate-dismissal} extends this to
near-singular systems.
The SVD of the KKT matrix provides the smallest singular
value~$\sigma_{\min}$ and the associated right singular
vector~$(\delta\beta, \delta\lambda, \delta\nu)$.
The three-valued version checks whether $\|\delta\beta\|$
is large enough to apply the proposition:
\[
  \tilde P_{\text{dismiss}} =
  \begin{cases}
    \textsc{true}
      & \text{if } \|\delta\beta\| > \varepsilon / \sigma_C
        \text{ (ensures } \delta\beta' \neq 0
        \text{ after projection; proposition applies),} \\
    \textsc{false}
      & \text{if } \sigma_{\min}
        \text{ is well-separated from zero (system is full-rank),} \\
    \textsc{unknown}
      & \text{if } \sigma_{\min} \text{ is small but }
        \|\delta\beta\| \leq \varepsilon / \sigma_C.
  \end{cases}
\]
When $\tilde P_{\text{dismiss}} = \textsc{unknown}$,
the near-singular direction is confined to the multiplier
block~$(\delta\lambda, \delta\nu)$:
the $\beta$-component of the solution is approximately unique,
and the algorithm uses it directly
(see Remark~\ref{rem:approximate-dismissal-algorithm}).

\begin{remark}[Parameter scales in practice]\label{rem:parameter-scales}
For polytopes generated by our random sampling
(Section~\ref{sec:random-sweep}),
the relevant quantities have the following typical magnitudes:
the symplectic inner products
$\omega_0(n_i, n_j) \in [-1, 1]$ (unit normals),
with $\omega_0(n_i, n_j) \approx 0$ occurring primarily
for near-Lagrangian 2-faces;
heights $h_i \in [h_{\min}, h_{\max}]$ with typical range
$[0.1, 10]$;
vertex norms $\|v\| \lesssim 100$;
and empirically $\beta_i = O(\|v\|)$, though the precise
scaling of $\beta$ with the polytope geometry (via the
constraints $N^\top \beta = 0$, $\eta^\top \beta = 1$)
deserves further investigation.
Before the normalisation $\eta^\top \beta = 1$,
the un-normalised period~$T$ of the orbit is bounded below
by the inradius (determined by $\min_k h_k$) and above by
the circumradius (determined by $\max_k \|v_k\|$);
the normalisation rescales $\beta$ by $1/T$.

These scales determine how aggressively the three-valued
predicates can distinguish $\textsc{true}$ from
$\textsc{unknown}$:
tolerances that are small relative to these scales yield
few $\textsc{unknown}$ verdicts.
Note that bounds on $h_k$ alone do not give good bounds on
vertex norms~$\|v_k\|$;
the random polytope generator should be checked to confirm
that extreme vertex norms are not generated.
\end{remark}
